\subsection{Standard Perturbation Theory (SPT) and Mode Coupling}
\label{subsec:spt_methodology}

Beyond the linear regime, non-linear mode coupling induces non-Gaussian signatures in the density field. 
In Eulerian Standard Perturbation Theory (SPT), the total density contrast up to second order is expanded as $\delta(\mathbf{x}) = \delta_1(\mathbf{x}) + \delta_2(\mathbf{x})$.

Given the primary linear field $\delta_1(\mathbf{k})$, the second-order field $\delta_2(\mathbf{x})$ is constructed in real space using FFT operations to evaluate advection and tidal shear interactions
    \begin{equation}
        \delta_2(\mathbf{x}) = \frac{17}{21} \delta_1^2(\mathbf{x}) + \mathbf{\Psi}^{(1)}(\mathbf{x}) \cdot \nabla \delta_1(\mathbf{x}) + \frac{2}{7} s_{ij}^2(\mathbf{x}),
        \label{eq:spt_delta2_real_space}
    \end{equation}
where $s_{ij} = \left(\nabla_i \nabla_j \nabla^{-2} - 1/3\delta_{ij}\right) \delta_1$ is the linear tidal shear tensor, and $\mathbf{\Psi}^{(1)} = -\nabla \nabla^{-2} \delta_1$ is the first-order displacement vector.

Non-Gaussian statistics induced by mode coupling are quantified via two primary observables
    \begin{enumerate}
    
    \item \textbf{Equilateral Bispectrum $B(k, k, k)$}: Measured from shell-filtered Fourier modes. At tree-level, the equilateral bispectrum target in Einstein--de Sitter (EdS) is given by
        \begin{equation}
            B_{\mathrm{tree}}^{\mathrm{equil}}(k) = \frac{12}{7} P_L^2(k).
            \label{eq:tree_level_bispectrum}
        \end{equation}
        
    \item \textbf{Smoothed Reduced Skewness $S_3(R)$}: Evaluated by smoothing the combined field $\delta(\mathbf{x})$ with a spherical Top-Hat window function $W_R(k) = 3(\sin kR - kR \cos kR)/(kR)^3$ of radius $R$
    \begin{equation}
        S_3(R) \equiv \frac{\langle \delta_R^3 \rangle}{\langle \delta_R^2 \rangle^2}.
        \label{eq:reduced_skewness}
    \end{equation}
\end{enumerate}